\section{Henneberg moves}
\label{sec:henneberg}

Henneberg's theorem characterises Laman graphs as the graphs obtained
from $K_2$ by iterating two combinatorial moves --- the \emph{Type~I}
and \emph{Type~II} extensions described below. This section formalises
the two moves, proves each preserves the Laman property, and develops
the \emph{iso half} of the structural decomposition: every Laman graph
on $n \ge 3$ vertices is isomorphic to a Type~I or Type~II extension
of \emph{some} graph $G'$ on $n - 1$ vertices. The deeper \emph{Laman
half} --- that some choice of $G'$ is itself Laman --- is the Phase~5
work covered in a later chapter.

Throughout, the fresh vertex introduced by either move is encoded as
$\mathrm{none} : \mathrm{Option}\,V$, with the old vertices embedded
via $\mathrm{some}$. This is mathlib's canonical ``add one element''
idiom and avoids a freshness side-condition on every move.

\subsection{Definitions}

\begin{definition}[Type~I move]
  \label{def:typeI}
  \lean{SimpleGraph.Henneberg.typeI}
  \leanok
  Given $G : \mathrm{SimpleGraph}\,V$ and two vertices $a, b \in V$, the
  \emph{Type~I extension} $\mathrm{typeI}\,G\,a\,b$ is the graph on
  $\mathrm{Option}\,V$ obtained by adding a fresh vertex
  $w = \mathrm{none}$ adjacent to (the images under $\mathrm{some}$ of)
  exactly $a$ and $b$. The other adjacencies are those of $G$.
  Under $a \ne b$ the fresh vertex has degree~2; if $a = b$ the parallel
  edges collapse to a single one and the move preserves the Laman
  property only under $a \ne b$.
\end{definition}

\begin{definition}[Type~II move]
  \label{def:typeII}
  \lean{SimpleGraph.Henneberg.typeII}
  \leanok
  Given $G : \mathrm{SimpleGraph}\,V$ and three vertices $a, b, c \in V$,
  the \emph{Type~II extension} $\mathrm{typeII}\,G\,a\,b\,c$ is the
  graph on $\mathrm{Option}\,V$ obtained by deleting the edge
  $\{a, b\}$ from $G$ and adding a fresh vertex $w = \mathrm{none}$
  adjacent to $a$, $b$, and $c$.

  The move is well-defined for arbitrary $a, b, c$; it preserves the
  Laman property only under the four hypotheses
  $a \ne b$, $c \ne a$, $c \ne b$, and $\{a, b\} \in E(G)$.
\end{definition}

\subsection{Edge-set decomposition and counts}

The two moves admit a clean edge-set decomposition into ``old'' edges
(image of $E(G)$ under $\mathrm{some}$, with one edge removed in the
Type~II case) plus the two-or-three new fresh edges incident to
$\mathrm{none}$. The cardinality consequences feed the tightness
half of the preservation theorems below.

\begin{lemma}[Edge set of Type~I]
  \label{lem:typeI-edgeSet}
  \lean{SimpleGraph.Henneberg.typeI_edgeSet}
  \leanok
  \uses{def:typeI}
  \[
    E(\mathrm{typeI}\,G\,a\,b)
    \;=\;
    \mathrm{Sym2.map\;some}\,[E(G)]
    \;\cup\;
    \bigl\{\,
      \{\mathrm{none}, \mathrm{some}\,a\},
      \;
      \{\mathrm{none}, \mathrm{some}\,b\}
    \,\bigr\}.
  \]
\end{lemma}
\begin{proof}
  \leanok
  By case analysis on the two endpoints of an edge in
  $\mathrm{Option}\,V$, the only edges incident to $\mathrm{none}$
  are those joining $\mathrm{none}$ to $\mathrm{some}\,a$ or
  $\mathrm{some}\,b$ (by definition of $\mathrm{typeI}$); all other
  edges have both endpoints in $\mathrm{some}\,[V]$ and correspond
  bijectively to edges of $G$.
\end{proof}

\begin{lemma}[Edge count of Type~I]
  \label{lem:typeI-edgeSet-ncard}
  \lean{SimpleGraph.Henneberg.typeI_edgeSet_ncard}
  \leanok
  \uses{lem:typeI-edgeSet}
  If $V$ is finite and $a \ne b$, then
  \[
    |E(\mathrm{typeI}\,G\,a\,b)| \;=\; |E(G)| + 2.
  \]
\end{lemma}
\begin{proof}
  \leanok
  Use \cref{lem:typeI-edgeSet}: the two summands are disjoint (each new
  edge contains $\mathrm{none}$, which is not in the image of
  $\mathrm{some}$), $\mathrm{Sym2.map\;some}$ is injective, and the
  two new edges are distinct because $a \ne b$.
\end{proof}

\begin{lemma}[Edge set of Type~II]
  \label{lem:typeII-edgeSet}
  \lean{SimpleGraph.Henneberg.typeII_edgeSet}
  \leanok
  \uses{def:typeII}
  \[
    E(\mathrm{typeII}\,G\,a\,b\,c)
    \;=\;
    \mathrm{Sym2.map\;some}\bigl[E(G) \setminus \{\{a,b\}\}\bigr]
    \;\cup\;
    \bigl\{\,
      \{\mathrm{none}, \mathrm{some}\,a\}, \;
      \{\mathrm{none}, \mathrm{some}\,b\}, \;
      \{\mathrm{none}, \mathrm{some}\,c\}
    \,\bigr\}.
  \]
\end{lemma}
\begin{proof}
  \leanok
  Case analysis as in \cref{lem:typeI-edgeSet}: edges incident to
  $\mathrm{none}$ are exactly the three new ones, edges with both
  endpoints in $\mathrm{some}\,[V]$ correspond bijectively to edges of
  $G$, and the $\{a, b\}$ edge is explicitly removed by the
  $s(u, v) \ne s(a, b)$ clause in the definition.
\end{proof}

\begin{lemma}[Edge count of Type~II]
  \label{lem:typeII-edgeSet-ncard}
  \lean{SimpleGraph.Henneberg.typeII_edgeSet_ncard}
  \leanok
  \uses{lem:typeII-edgeSet}
  If $V$ is finite and $a \ne b$, $c \ne a$, $c \ne b$, then
  \[
    |E(\mathrm{typeII}\,G\,a\,b\,c)|
    \;=\;
    |E(G) \setminus \{\{a,b\}\}| + 3.
  \]
  The deletion is recorded on the right-hand side rather than
  consumed by a $\{a, b\} \in E(G)$ hypothesis, so the lemma does not
  require $G$-adjacency of $a$ and $b$.
\end{lemma}
\begin{proof}
  \leanok
  As in \cref{lem:typeI-edgeSet-ncard}, the two summands of
  \cref{lem:typeII-edgeSet} are disjoint and the
  $\mathrm{Sym2.map\;some}$ side is injective; the three new edges
  are pairwise distinct because $a, b, c$ are.
\end{proof}

\subsection{Preservation of the Laman property}

\begin{theorem}[Type~I preserves Laman]
  \label{thm:typeI-isLaman}
  \lean{SimpleGraph.Henneberg.typeI_isLaman}
  \leanok
  \uses{def:typeI, def:isLaman, lem:typeI-edgeSet-ncard}
  If $V$ is finite, $G$ is Laman, and $a \ne b$, then
  $\mathrm{typeI}\,G\,a\,b$ is Laman.
\end{theorem}
\begin{proof}
  \leanok
  Tightness follows from \cref{lem:typeI-edgeSet-ncard} and
  $|\mathrm{Option}\,V| = |V| + 1$: starting from
  $|E(G)| + 3 = 2|V|$, both sides gain $+2$.

  For sparsity, fix $s \subseteq \mathrm{Option}\,V$ with
  $3 \le 2|s|$ and write $s' := s \cap \mathrm{some}[V]$ for the
  $\mathrm{some}$-preimage. The edge set
  $\edgesIn[\mathrm{typeI}\,G\,a\,b]{s}$ decomposes as the $G$-edges
  on $s'$ plus the subset $T$ of the two new edges
  $\{\mathrm{none}, \mathrm{some}\,a\}$,
  $\{\mathrm{none}, \mathrm{some}\,b\}$ that lie in $s$. Case-split:
  \begin{itemize}
    \item If $\mathrm{none} \notin s$, then $T = \emptyset$ and
      $|s| = |s'|$, so the sparsity of $G$ at $s'$ transfers directly.
    \item If $\mathrm{none} \in s$, then $|s| = |s'| + 1$ and
      $|T| \le 2$. Sub-case on $|s'|$: if $|s'| = 0$ the precondition
      fails; if $|s'| = 1$, say $s' = \{w\}$, then no $G$-edge lies
      in $s'$ and $T \subseteq \{\{\mathrm{none}, \mathrm{some}\,w\}\}$
      so $|T| \le 1$, matching $2|s| - 3 = 0$; if $|s'| \ge 2$, apply
      sparsity of $G$ at $s'$ and use $|T| \le 2$.
  \end{itemize}
\end{proof}

\begin{theorem}[Type~II preserves Laman]
  \label{thm:typeII-isLaman}
  \lean{SimpleGraph.Henneberg.typeII_isLaman}
  \leanok
  \uses{def:typeII, def:isLaman, lem:typeII-edgeSet-ncard}
  If $V$ is finite, $G$ is Laman, the three vertices $a, b, c$ are
  pairwise distinct, and $\{a, b\} \in E(G)$, then
  $\mathrm{typeII}\,G\,a\,b\,c$ is Laman.
\end{theorem}
\begin{proof}
  \leanok
  Tightness: by \cref{lem:typeII-edgeSet-ncard} and
  $\{a, b\} \in E(G)$, the deletion drops the count by exactly $1$ and
  three new edges arrive, so $|E| + 3 = 2|V|$ shifts to
  $|E| + 5 = 2(|V| + 1) = 2|V| + 2$, which simplifies to the new
  tightness identity.

  Sparsity mirrors \cref{thm:typeI-isLaman}. The new wrinkle is that
  the deleted edge $\{a, b\}$ interacts with the (up to) three new
  fresh edges. Uniform $+2$ bound on the net gain: in the
  $\mathrm{none} \in s$, $|s'| \ge 2$ subcase, either
  $\{a, b\} \in \edgesIn{s'}$ (deletion loses $1$, the three new
  edges contribute $\le 3$, net $\le 2$) or $\{a, b\} \notin \edgesIn{s'}$
  (deletion loses $0$; since $\{a, b\} \in E(G)$ this forces
  $a \notin s'$ or $b \notin s'$, so the corresponding new edge
  $\{\mathrm{none}, \mathrm{some}\,a\}$ or
  $\{\mathrm{none}, \mathrm{some}\,b\}$ is also absent, bounding
  $|T'| \le 2$). The remaining sub-cases ($|s'| \in \{0, 1\}$ and
  $\mathrm{none} \notin s$) match \cref{thm:typeI-isLaman}.
\end{proof}

\subsection{Decomposition: iso half}
\label{sec:henneberg-decomp}

A Laman graph on $n \ge 3$ vertices is, up to canonical relabelling,
the result of either a Type~I or a Type~II move applied to \emph{some}
graph $G'$ on $\{w \in V : w \ne v\}$, where $v$ is a chosen vertex of
degree $2$ or $3$. The statement below does \textbf{not} assert
$G'$ is itself Laman --- that is the deeper Phase~5 direction and is
known to fail for arbitrary choices of bridging pair in the Type~II
case.

\begin{theorem}[Iso decomposition]
  \label{thm:isLaman-exists-typeI-or-typeII-iso}
  \lean{SimpleGraph.IsLaman.exists_typeI_or_typeII_iso}
  \leanok
  \uses{def:isLaman, def:typeI, def:typeII,
    lem:isLaman-exists-two-le-degree-le-three,
    lem:isLaman-exists-nonadj-among-three-neighbors}
  Let $G$ be Laman on a finite vertex type $V$ with $|V| \ge 3$. Then
  there exist $v \in V$ and a graph $G'$ on $\{w : V \mid w \ne v\}$
  such that either
  \begin{enumerate}
    \item there exist $a, b$ with $a \ne b$ and a graph isomorphism
      $G \simeq \mathrm{typeI}\,G'\,a\,b$; or
    \item there exist $a, b, c$ with $a \ne b$, $c \ne a$, $c \ne b$,
      $\{a, b\} \in E(G')$, and a graph isomorphism
      $G \simeq \mathrm{typeII}\,G'\,a\,b\,c$.
  \end{enumerate}
\end{theorem}
\begin{proof}
  \leanok
  By \cref{lem:isLaman-exists-two-le-degree-le-three} pick $v$ with
  $2 \le \deg_G(v) \le 3$, and let $G'$ be the induced subgraph
  $G[\{w \ne v\}]$ (augmented for the Type~II case below). The
  underlying type-level equivalence is
  $V \simeq \mathrm{Option}\,\{w : V \mid w \ne v\}$, mapping $v$ to
  $\mathrm{none}$.

  \emph{Degree~$2$.} Let $\{a, b\} = N_G(v)$; then
  $N_G(v) = \{a, b\}$ and $a \ne b$, so the type-level equivalence
  upgrades to a graph iso $G \simeq \mathrm{typeI}\,G'\,a\,b$.

  \emph{Degree~$3$.} Let $\{a, b, c\} = N_G(v)$. By
  \cref{lem:isLaman-exists-nonadj-among-three-neighbors} at least one
  pair among $\{a, b\}, \{a, c\}, \{b, c\}$ is non-adjacent in $G$;
  relabel so this is $\{a, b\}$. Let $G'$ be the induced subgraph
  $G[\{w \ne v\}]$ \emph{augmented} with the bridging edge
  $\{a, b\}$. The type-level equivalence again gives a graph iso
  $G \simeq \mathrm{typeII}\,G'\,a\,b\,c$: the bridging edge in $G'$
  is exactly the one the Type~II definition removes.
\end{proof}

\subsection{Worked example: \texorpdfstring{$K_4 \setminus e$}{K4 minus e}}

\begin{theorem}[$K_4 \setminus e$ is Laman]
  \label{thm:top-fin-four-minus-edge-isLaman}
  \lean{SimpleGraph.top_fin_four_minus_edge_isLaman}
  \leanok
  \uses{thm:typeI-isLaman, thm:top-fin-two-isLaman, lem:isLaman-iso}
  The graph $K_4$ on the vertex set $\mathrm{Fin}\,4$, minus the single
  edge $\{2, 3\}$, is Laman.
\end{theorem}
\begin{proof}
  \leanok
  Apply \cref{thm:typeI-isLaman} twice from $K_2$
  (\cref{thm:top-fin-two-isLaman}): a single Type~I move yields a
  Laman graph on three vertices ($K_3$ in disguise), and a second
  Type~I move yields a Laman graph on four vertices. This Laman graph
  on $\mathrm{Option}\,(\mathrm{Option}\,(\mathrm{Fin}\,2))$ is iso to
  $K_4 \setminus \{2, 3\}$ on $\mathrm{Fin}\,4$ via an explicit
  bijection; transport along \cref{lem:isLaman-iso}.
\end{proof}
